\documentclass[11pt,a4paper]{article}

\usepackage[spanish,es-tabla]{babel}
\usepackage[utf8]{inputenc}
\usepackage[T1]{fontenc}
\usepackage{graphicx}
\usepackage{geometry}
\usepackage{fancyhdr}
\usepackage{tcolorbox}
\tcbuselibrary{skins,breakable}
\usepackage{xcolor}
\usepackage{hyperref}
\usepackage{fontawesome5}
\usepackage{enumitem}
\usepackage{titlesec}
\usepackage{setspace}

\geometry{
 a4paper,
 left=1.8cm,
 right=1.8cm,
 top=2.2cm,
 bottom=2.2cm
}

% --- Definición de Colores ---
\definecolor{puceBlue}{HTML}{003366}
\definecolor{huecaOrange}{HTML}{F47B20}
\definecolor{huecaOrangeDark}{HTML}{C85A00}
\definecolor{adminNavy}{HTML}{101828}
\definecolor{adminGold}{HTML}{FFB020}
\definecolor{cardBg}{HTML}{F8FAFC}
\definecolor{cardBorder}{HTML}{E2E8F0}
\definecolor{tagGreen}{HTML}{16A34A}
\definecolor{textPrimary}{HTML}{1E293B}
\definecolor{textSecondary}{HTML}{64748B}

\hypersetup{
  colorlinks=true,
  linkcolor=puceBlue,
  urlcolor=puceBlue,
  citecolor=puceBlue
}

% --- Estilos de Encabezado y Pie ---
\pagestyle{fancy}
\fancyhf{}
\fancyhead[L]{\small\textbf{Hue-Quito} --- Catálogo Visual de Interfaces Móviles}
\fancyhead[R]{\small PUCE 2026}
\fancyfoot[L]{\small\textcolor{textSecondary}{Proyecto de Titulación --- L. Rodríguez \& J. Guaña}}
\fancyfoot[R]{\small\textbf{\thepage}}
\renewcommand{\headrulewidth}{0.4pt}
\renewcommand{\footrulewidth}{0.4pt}

% --- Macro para Separador de Sección ---
\newcommand{\sectionheader}[3]{%
  \newpage
  \vspace*{2cm}
  \begin{center}
    \begin{tcolorbox}[
      enhanced,
      colback=#2,
      colframe=#2,
      arc=4mm,
      width=0.92\textwidth,
      halign=center,
      valign=center,
      top=1.5cm,
      bottom=1.5cm
    ]
      {\Huge\bfseries\color{white} #1}\\[0.6cm]
      {\large\color{white!90} #3}
    \end{tcolorbox}
  \end{center}
  \vfill
  \newpage
}

% --- Macro para Pantalla Individual ---
\newcommand{\pantallashowcase}[7]{%
  % #1: Número y Título de la Pantalla
  % #2: Archivo fuente .dart
  % #3: Rol Badge Text
  % #4: Rol Badge Color
  % #5: Ruta de la imagen
  % #6: Altura de la imagen (ej. 17.5cm)
  % #7: Descripción y detalles estructurados
  \newpage
  \begin{tcolorbox}[
    enhanced,
    colback=white,
    colframe=cardBorder,
    arc=2.5mm,
    boxrule=0.8pt,
    left=4mm, right=4mm, top=3mm, bottom=3mm
  ]
    \begin{minipage}[c]{0.70\textwidth}
      {\large\bfseries\color{textPrimary} #1}\\[0.1cm]
      {\footnotesize\ttfamily\color{textSecondary}\faIcon{file-code}\ \nolinkurl{#2}}
    \end{minipage}%
    \hfill
    \begin{minipage}[c]{0.28\textwidth}
      \flushright
      \tcbox[
        enhanced,
        colback=#4!10,
        colframe=#4,
        boxrule=0.6pt,
        arc=1.5mm,
        top=1mm, bottom=1mm, left=2.5mm, right=2.5mm
      ]{\footnotesize\bfseries\color{#4} #3}
    \end{minipage}
  \end{tcolorbox}
  \vspace{0.2cm}

  \noindent
  \begin{minipage}[t]{0.45\textwidth}
    \centering
    \vspace{0pt}
    \begin{tcolorbox}[
      enhanced,
      colback=white,
      colframe=gray!30,
      arc=3mm,
      boxrule=0.7pt,
      halign=center,
      top=2mm, bottom=2mm, left=2mm, right=2mm
    ]
      \includegraphics[height=#6, width=\linewidth, keepaspectratio]{#5}
    \end{tcolorbox}
  \end{minipage}%
  \hfill
  \begin{minipage}[t]{0.52\textwidth}
    \vspace{0pt}
    #7
  \end{minipage}
}

% Tras el título de un módulo, la primera pantalla va en la misma página
\newif\ifskipnewpage

% --- Captura que aún no existe: se muestra un aviso en lugar de romper la compilación ---
\newcommand{\capturaimg}[2]{%
  \IfFileExists{#1}{%
    \includegraphics[height=#2, width=\linewidth, keepaspectratio]{#1}%
  }{%
    \begin{minipage}[c][#2][c]{\linewidth}\centering\color{textSecondary}\small\faIcon{camera}\ Captura pendiente\end{minipage}%
  }%
}

% --- Pantalla con DOS capturas lado a lado y la descripción debajo ---
\newcommand{\pantalladoble}[8]{%
  % #1: Título  #2: archivo .dart  #3: rol  #4: color del rol
  % #5: captura izquierda  #6: captura derecha  #7: altura  #8: descripción
  \ifskipnewpage\global\skipnewpagefalse\else\newpage\fi
  \begin{tcolorbox}[
    enhanced,
    colback=white,
    colframe=cardBorder,
    arc=2.5mm,
    boxrule=0.8pt,
    left=4mm, right=4mm, top=3mm, bottom=3mm
  ]
    \begin{minipage}[c]{0.70\textwidth}
      {\large\bfseries\color{textPrimary} #1}\\[0.1cm]
      {\footnotesize\ttfamily\color{textSecondary}\faIcon{file-code}\ \nolinkurl{#2}}
    \end{minipage}%
    \hfill
    \begin{minipage}[c]{0.28\textwidth}
      \flushright
      \tcbox[
        enhanced,
        colback=#4!10,
        colframe=#4,
        boxrule=0.6pt,
        arc=1.5mm,
        top=1mm, bottom=1mm, left=2.5mm, right=2.5mm
      ]{\footnotesize\bfseries\color{#4} #3}
    \end{minipage}
  \end{tcolorbox}
  \vspace{0.2cm}

  \noindent
  \begin{minipage}[t]{0.49\textwidth}
    \begin{tcolorbox}[enhanced, colback=white, colframe=gray!30, arc=3mm, boxrule=0.7pt, halign=center, top=2mm, bottom=2mm, left=2mm, right=2mm]
      \capturaimg{#5}{#7}
    \end{tcolorbox}
  \end{minipage}%
  \hfill
  \begin{minipage}[t]{0.49\textwidth}
    \begin{tcolorbox}[enhanced, colback=white, colframe=gray!30, arc=3mm, boxrule=0.7pt, halign=center, top=2mm, bottom=2mm, left=2mm, right=2mm]
      \capturaimg{#6}{#7}
    \end{tcolorbox}
  \end{minipage}

  \vspace{0.2cm}
  #8
}

% La macro de una sola captura usa también el aviso de captura pendiente
\renewcommand{\pantallashowcase}[7]{%
  \ifskipnewpage\global\skipnewpagefalse\else\newpage\fi
  \begin{tcolorbox}[
    enhanced,
    colback=white,
    colframe=cardBorder,
    arc=2.5mm,
    boxrule=0.8pt,
    left=4mm, right=4mm, top=3mm, bottom=3mm
  ]
    \begin{minipage}[c]{0.70\textwidth}
      {\large\bfseries\color{textPrimary} #1}\\[0.1cm]
      {\footnotesize\ttfamily\color{textSecondary}\faIcon{file-code}\ \nolinkurl{#2}}
    \end{minipage}%
    \hfill
    \begin{minipage}[c]{0.28\textwidth}
      \flushright
      \tcbox[
        enhanced,
        colback=#4!10,
        colframe=#4,
        boxrule=0.6pt,
        arc=1.5mm,
        top=1mm, bottom=1mm, left=2.5mm, right=2.5mm
      ]{\footnotesize\bfseries\color{#4} #3}
    \end{minipage}
  \end{tcolorbox}
  \vspace{0.2cm}

  \noindent
  \begin{minipage}[t]{0.45\textwidth}
    \centering
    \vspace{0pt}
    \begin{tcolorbox}[enhanced, colback=white, colframe=gray!30, arc=3mm, boxrule=0.7pt, halign=center, top=2mm, bottom=2mm, left=2mm, right=2mm]
      \capturaimg{#5}{#6}
    \end{tcolorbox}
  \end{minipage}%
  \hfill
  \begin{minipage}[t]{0.52\textwidth}
    \vspace{0pt}
    #7
  \end{minipage}
}

\newcommand{\proposito}[1]{%
  \begin{tcolorbox}[enhanced, colback=cardBg, colframe=cardBorder, arc=2mm, boxrule=0.6pt]
    \textbf{\large Propósito}\\[0.1cm]
    #1
  \end{tcolorbox}
  \vspace{0.2cm}
}

\begin{document}

% =============================================================================
% PORTADA INSTITUCIONAL PUCE
% =============================================================================
\begin{titlepage}
    \centering
    \vspace*{0.5cm}
    {\large \textbf{PONTIFICIA UNIVERSIDAD CATÓLICA DEL ECUADOR}}\\[0.6cm]
    {\Large \textbf{PROYECTO DE TITULACIÓN}}\\[1.2cm]
    {\Large \textbf{TÍTULO DEL PROYECTO:}}\\[0.4cm]
    {\LARGE\bfseries Plataforma Turística Innovadora con Gamificación,\\Inteligencia Artificial y Logística Inteligente\par}
    \vspace{0.6cm}
    {\large\textbf{Aplicación Móvil: Hue-Quito (Descubre las Huecas de Quito)}\par}
    \vspace{0.3cm}
    {\normalsize\textbf{\color{huecaOrangeDark} ANEXO DE EVIDENCIAS: CATÁLOGO VISUAL DE INTERFACES}\par}
    \vspace{0.2cm}
    {\small\color{textSecondary} Versión 2 --- actualizada con zonas de Quito, menú con buscador y tours personalizados\par}
    \vspace{1.0cm}

    % LOGO PUCE
    \IfFileExists{logo-puce.png}{%
      \includegraphics[width=0.28\textwidth]{logo-puce.png}\\[1cm]
    }{%
      \IfFileExists{docs/logo-puce.png}{%
        \includegraphics[width=0.28\textwidth]{docs/logo-puce.png}\\[1cm]
      }{%
        \vspace{2.5cm}
      }%
    }
    \vspace{0.8cm}

    {\large \textbf{AUTOR:}}\\[0.2cm]
    {\large LUIS RODRÍGUEZ}\\[1cm]
    {\large \textbf{DOCENTE DIRECTOR:}}\\[0.2cm]
    {\large JAVIER GUAÑA}\\[1.5cm]

    \vfill
    {\large Quito -- Ecuador}\\
    {\large 2026}
\end{titlepage}

% =============================================================================
% RESUMEN Y CONTENIDO
% =============================================================================
\newpage
\section*{Presentación del Catálogo Visual}

El presente documento constituye el \textbf{Anexo Visual y Funcional} del Proyecto de Titulación \textit{``Plataforma Turística Innovadora con Gamificación, Inteligencia Artificial y Logística Inteligente (Hue-Quito)''}.

Reúne capturas reales de pantalla de la versión vigente de la aplicación, tomadas en formato de teléfono móvil sobre la versión publicada (\url{https://huequito-b091a.web.app}), que comparte el mismo código fuente que la aplicación Android. Los datos visibles corresponden al catálogo real de \textbf{20 huecas tradicionales de Quito}, organizadas en cuatro zonas: \textit{Centro, Norte, Sur y Valles}. Las pantallas del dueño y del administrador se capturaron sobre una copia local de la aplicación conectada a los emuladores de Firebase, con las mismas 20 huecas y cuentas de prueba (por ejemplo, la dueña ficticia de \textit{El Pretil}), para no alterar los datos reales. Las interfaces se han clasificado en los tres roles que operan dentro de la plataforma:

\begin{enumerate}[label=\textbf{\Roman*.}]
  \item \textbf{Módulo de Usuario (Comensal / Turista):} Cuestionario de perfil gastronómico por zona, exploración de huecas, menú con buscador y filtros por categoría, tours personalizados optimizados por cercanía, navegación paso a paso, escaneo QR de sellos, álbum de cromos por zona y preferencias de accesibilidad.
  \item \textbf{Módulo de Negocio (Dueño / Encargado de Hueca):} Tablero de métricas en tiempo real, menú digital organizado por categorías con ayuda de Inteligencia Artificial (Gemini), emisor de QR antifraude, redactor publicitario con IA y gestión de sucursales.
  \item \textbf{Módulo de Administración (Superadministrador):} Centro de control global, auditoría de huecas, zonas y cromos, moderación de avisos, revisión de solicitudes de nuevas huecas y gestión de cuentas comerciales.
\end{enumerate}

\vspace{0.5cm}
\begin{tcolorbox}[
  enhanced,
  colback=huecaOrange!5,
  colframe=huecaOrange,
  arc=2.5mm,
  boxrule=0.8pt,
  title=\textbf{\faIcon{info-circle}\ Nota de Integridad Técnica}
]
  Todas las interfaces presentadas corresponden al código fuente alojado en el repositorio oficial del proyecto, implementadas en el framework \textbf{Flutter} con gestión de estado reactiva (\textbf{Riverpod}), backend en tiempo real (\textbf{Cloud Firestore}), autenticación federada (\textbf{Firebase Auth}), mapas y rutas con \textbf{Google Maps Platform}, lectura de cartas con \textbf{Gemini} a través de \textbf{Firebase AI Logic} (sin claves de API dentro de la aplicación) y reconocimiento de texto en el dispositivo (\textbf{Google ML Kit}).
\end{tcolorbox}

\vspace{0.8cm}
\tableofcontents
\newpage

% =============================================================================
% PARTE I: MÓDULO DE USUARIO (COMENSAL)
% =============================================================================
\section{Módulo de Usuario (Comensal / Turista)}
\skipnewpagetrue

% --- PANTALLA 1: WelcomeScreen ---
\pantallashowcase
  {Pantalla 1: Bienvenida a Hue-Quito}
  {lib/screens/welcome_screen.dart}
  {COMENSAL}
  {huecaOrange}
  {capturas/01_bienvenida.png}
  {17.5cm}
  {
    \proposito{Primera pantalla de la aplicación. Presenta la propuesta de valor y permite empezar el cuestionario de perfil o ir directo al inicio de sesión.}

    \textbf{Componentes Clave:}
    \begin{itemize}[leftmargin=*, itemsep=2pt]
      \item Logotipo de Hue-Quito sobre fotografía de ambiente de una hueca tradicional.
      \item Indicador de ciudad (\textit{Quito · Ecuador}) y selector de idioma (Español / English).
      \item Tres tarjetas con los pilares de la app: \textit{Huecas Tradicionales}, \textit{Mapa \& Circuitos Guiados} y \textit{Cartilla de Sellos \& Cromos}.
      \item Botón principal \textit{``Comenzar Aventura''} que abre el cuestionario.
      \item Accesos secundarios: \textit{Saltar cuestionario e iniciar sesión} e \textit{Iniciar Sesión}.
    \end{itemize}
  }

% --- PANTALLA 2: PreferencesScreen (zona + pregunta) ---
\pantalladoble
  {Pantalla 2: Cuestionario de Perfil Gastronómico}
  {lib/screens/preferences_screen.dart}
  {COMENSAL}
  {huecaOrange}
  {capturas/02_quiz_zonas.png}
  {capturas/03_quiz_pregunta.png}
  {14cm}
  {
    \proposito{Cuestionario inicial de 5 pasos que define la zona del usuario y su perfil de antojos para personalizar las recomendaciones.}
    \begin{itemize}[leftmargin=*, itemsep=2pt]
      \item \textbf{Paso 1 --- Tu Sector (izquierda):} Quito se divide en solo cuatro zonas: \textit{Centro, Norte, Sur y Valles}. Cada tarjeta muestra los barrios que incluye (por ejemplo, Norte agrupa La Mariscal, La Floresta, La Carolina e Iñaquito) y cuántas huecas hay en ella. La opción \textit{Otro sector} permite explorar toda la ciudad.
      \item \textbf{Pasos 2 a 5 --- Preguntas de antojo (derecha):} Cuatro preguntas con fotografía y respuestas ilustradas; las preguntas son editables por el administrador.
      \item Barra de progreso, selector de idioma y botón de regreso en cada paso.
    \end{itemize}
  }

% --- PANTALLA 3: Resultado del quiz ---
\pantallashowcase
  {Pantalla 3: Resultado del Perfil Gastronómico}
  {lib/screens/preferences_screen.dart}
  {COMENSAL}
  {huecaOrange}
  {capturas/04_quiz_resultado.png}
  {17.5cm}
  {
    \proposito{Resume el perfil culinario del usuario calculado a partir de sus respuestas.}

    \textbf{Componentes Clave:}
    \begin{itemize}[leftmargin=*, itemsep=2pt]
      \item Perfil asignado con ilustración y descripción (ejemplo: \textit{El Aventurero}).
      \item Insignia con la zona elegida (\textit{Tu Zona: Centro}).
      \item Etiquetas de afinidad usadas por el algoritmo de recomendación (\textit{\#tripa mishqui, \#cuy, \#guatita, \#mercado}\dots).
      \item Acciones: \textit{Iniciar Sesión / Registrarme}, \textit{Continuar como invitado} y \textit{Repetir Quiz}.
    \end{itemize}
    \vspace{0.2cm}

    \textbf{Lógica y Datos:}
    \begin{itemize}[leftmargin=*, itemsep=2pt]
      \item El perfil se guarda en la cuenta del usuario o, si es invitado, en el propio celular.
    \end{itemize}
  }

% --- PANTALLA 4: LoginScreen ---
\pantallashowcase
  {Pantalla 4: Inicio de Sesión y Acceso}
  {lib/screens/login_screen.dart}
  {COMENSAL}
  {huecaOrange}
  {capturas/05_login.png}
  {17.5cm}
  {
    \proposito{Punto de entrada seguro a la plataforma que autentica al comensal o turista, permitiendo guardar su pasaporte gastronómico y sellos en la nube.}

    \textbf{Componentes Clave:}
    \begin{itemize}[leftmargin=*, itemsep=2pt]
      \item Logotipo e identidad visual de Hue-Quito.
      \item Campos de correo electrónico y contraseña con botón para mostrar u ocultar la contraseña.
      \item Enlace \textit{¿Olvidaste tu contraseña?} que envía un correo de recuperación.
      \item Inicio rápido con cuenta de Google.
      \item Opción \textit{``Continuar como Invitado''} para explorar sin registro.
      \item Enlace de registro de nuevas cuentas.
    \end{itemize}
    \vspace{0.2cm}

    \textbf{Lógica y Seguridad:}
    \begin{itemize}[leftmargin=*, itemsep=2pt]
      \item Integrado con \texttt{FirebaseAuth}; validación de formato de correo y longitud de contraseña antes de enviar la petición.
    \end{itemize}
  }

% --- PANTALLA 5: Guía del comensal ---
\pantallashowcase
  {Pantalla 5: Guía Interactiva del Comensal}
  {lib/widgets/client_onboarding_dialog.dart}
  {COMENSAL}
  {huecaOrange}
  {capturas/06_guia_comensal.png}
  {17.5cm}
  {
    \proposito{Tutorial de 6 pasos que aparece la primera vez que el usuario entra al inicio y explica cómo aprovechar la aplicación.}

    \textbf{Componentes Clave:}
    \begin{itemize}[leftmargin=*, itemsep=2pt]
      \item Tarjetas ilustradas con los temas: descubrir huecas, filtrar por zona, menús con precios, búsqueda por plato, rutas, sellos y cromos.
      \item Consejos prácticos destacados (\textit{Tip}).
      \item Indicador de avance, botón \textit{Siguiente} y opción \textit{Saltar}.
      \item Puede volver a abrirse desde el perfil en \textit{Tutorial de la App}.
    \end{itemize}
  }

% --- PANTALLA 6: HomeScreen ---
\pantalladoble
  {Pantalla 6: Exploración Principal y Recomendaciones}
  {lib/screens/home_screen.dart}
  {COMENSAL}
  {huecaOrange}
  {capturas/07_home.png}
  {capturas/08_home_zonas.png}
  {14cm}
  {
    \proposito{Pantalla de inicio que presenta las huecas recomendadas según el perfil y la zona del usuario.}
    \begin{itemize}[leftmargin=*, itemsep=2pt]
      \item \textbf{Selector de zona (derecha):} \textit{Todo Quito, Centro, Norte, Sur y Valles}. Las huecas con sectores antiguos o barrios (por ejemplo \textit{Conocoto} o \textit{La Mariscal}) se agrupan automáticamente en su zona.
      \item Buscador en tiempo real por nombre de hueca, plato o producto.
      \item Carrusel \textit{``Nuestra recomendación para ti''} con plato insignia, distancia, precio de referencia, calificación y etiquetas.
      \item Accesos rápidos \textit{Cómo llegar} (abre Google Maps o Waze) y \textit{Ver Hueca}.
      \item Categorías \textit{¿Qué se te antoja hoy?}: Platos Fuertes, Sopas, Mariscos, Dulces\dots
      \item Barra inferior: \textit{Explorar, Rutas, Escanear, Mi Álbum y Perfil}.
    \end{itemize}
  }

% --- PANTALLA 7: HuecaDetailScreen ---
\pantallashowcase
  {Pantalla 7: Ficha Detallada de la Hueca}
  {lib/screens/hueca_detail_screen.dart}
  {COMENSAL}
  {huecaOrange}
  {capturas/09_detalle.png}
  {17.5cm}
  {
    \proposito{Ficha integral del establecimiento tradicional (ejemplo: \textit{Corvinas de María Paz}, en el Mercado Central, zona Centro) con su información culinaria, logística y de contacto.}

    \textbf{Componentes Clave:}
    \begin{itemize}[leftmargin=*, itemsep=2pt]
      \item Portada con plato estrella, precio, dirección y visor de fotos a pantalla completa.
      \item Indicador de reputación (\textit{Nuevo} mientras no tenga reseñas) y zona.
      \item Botón de suscripción a avisos del local (campana).
      \item Botón destacado \textit{``Cómo llegar''}: abre la ubicación en Google Maps o Waze; si hay varias sucursales, primero pregunta a cuál ir.
      \item Acciones rápidas: \textit{Llamar}, \textit{Avisos} y \textit{Compartir}.
      \item Pestañas: \textit{Sobre la Hueca}, \textit{Menú Tradicional} y \textit{Ubicación}.
    \end{itemize}
  }

% --- PANTALLA 8: Menú con buscador ---
\pantalladoble
  {Pantalla 8: Menú con Buscador y Filtros por Categoría}
  {lib/screens/hueca_detail_screen.dart}
  {COMENSAL}
  {huecaOrange}
  {capturas/12_menu_buscador.png}
  {capturas/13_menu_filtrado.png}
  {14cm}
  {
    \proposito{Evita que los platos del final de un menú largo se pierdan: el comensal puede buscar y filtrar el menú de la hueca.}
    \begin{itemize}[leftmargin=*, itemsep=2pt]
      \item \textbf{Lupa de búsqueda:} aparece cuando el menú tiene más de 4 platos y filtra por nombre o descripción mientras se escribe (derecha: búsqueda \textit{``humita''}).
      \item \textbf{Chips de categoría:} \textit{Todo (8), Desayunos (1), Entradas (2)\dots}, con el número de platos de cada una.
      \item Platos agrupados bajo encabezados con ícono por categoría (Desayunos, Entradas, Platos Fuertes, Postres, Bebidas\dots).
      \item Cada plato muestra nombre, precio con dos decimales y foto.
      \item Las categorías las define el dueño de la hueca (ver Módulo de Negocio).
    \end{itemize}
  }

% --- PANTALLA 9: Mis Tours ---
\pantalladoble
  {Pantalla 9: Mis Tours Gastronómicos}
  {lib/screens/route_list_screen.dart}
  {COMENSAL}
  {huecaOrange}
  {capturas/15_rutas.png}
  {capturas/16_tour_asistente.png}
  {14cm}
  {
    \proposito{Cada usuario arma sus propios recorridos gastronómicos por Quito.}
    \begin{itemize}[leftmargin=*, itemsep=2pt]
      \item \textbf{Lista de tours guardados (izquierda):} tarjeta con número de paradas, distancia total y tiempo estimado en auto, paradas numeradas con su plato y zona, y botones \textit{Ver Tour} e \textit{Iniciar Ruta}.
      \item \textbf{Asistente (derecha):} el usuario elige el tamaño del recorrido: \textit{2 Huecas (Ruta Express)}, \textit{3 Huecas (Ruta Clásica)} o \textit{4 Huecas (Tour Glotón)}.
      \item Dos modos de creación: \textit{Tour Asistido por Preguntas} o \textit{Armar a mi Manera} desde el catálogo.
      \item Selector entre \textit{Mis Tours} y \textit{Mapa Interactivo}.
    \end{itemize}
  }

% --- PANTALLA 10: Asistente de paradas ---
\pantalladoble
  {Pantalla 10: Asistente de Paradas por Cercanía}
  {lib/screens/tour_wizard_screen.dart}
  {COMENSAL}
  {huecaOrange}
  {capturas/17_tour_parada.png}
  {capturas/18_tour_cercania.png}
  {14cm}
  {
    \proposito{Construye el tour parada por parada preguntando por el antojo del momento.}
    \begin{itemize}[leftmargin=*, itemsep=2pt]
      \item \textbf{Parada 1 (izquierda):} pregunta \textit{¿Cómo arrancamos esta aventura?} con chips de categoría (Sopas y Caldos, Platos Fuertes, Mariscos\dots) y huecas con su zona, calificación y plato.
      \item \textbf{Paradas siguientes (derecha):} las huecas se ordenan por cercanía a la parada anterior, con distancia, minutos en auto y la etiqueta \textit{Muy cerca}.
      \item Botón inferior para confirmar cada parada y continuar.
    \end{itemize}
  }

% --- PANTALLA 11: Resumen del tour ---
\pantallashowcase
  {Pantalla 11: Resumen del Tour Personalizado}
  {lib/screens/tour_wizard_screen.dart}
  {COMENSAL}
  {huecaOrange}
  {capturas/19_tour_resumen.png}
  {17.5cm}
  {
    \proposito{Muestra el recorrido final antes de guardarlo.}

    \textbf{Componentes Clave:}
    \begin{itemize}[leftmargin=*, itemsep=2pt]
      \item Mensaje \textit{¡Tu tour gastronómico está listo!}
      \item Campo editable con el nombre del tour.
      \item Métricas: número de huecas, distancia total y tiempo estimado en auto.
      \item Itinerario numerado con la distancia y el tiempo entre cada parada.
      \item Botones \textit{Guardar Tour en mi Perfil} y \textit{Guardar e Iniciar Recorrido}.
    \end{itemize}
    \vspace{0.2cm}

    \textbf{Lógica y Datos:}
    \begin{itemize}[leftmargin=*, itemsep=2pt]
      \item Con sesión iniciada el tour se guarda en Firestore; como invitado se guarda en el celular.
    \end{itemize}
  }

% --- PANTALLA 12: Navegación del tour ---
\pantallashowcase
  {Pantalla 12: Navegación Paso a Paso del Tour}
  {lib/screens/route_navigation_screen.dart}
  {COMENSAL}
  {huecaOrange}
  {capturas/27_navegacion.png}
  {17.5cm}
  {
    \proposito{Guía al comensal por todas las paradas del tour con la ruta trazada sobre el mapa.}

    \textbf{Componentes Clave:}
    \begin{itemize}[leftmargin=*, itemsep=2pt]
      \item Modo de transporte: \textit{A pie} o \textit{En carro}.
      \item Mapa con la ruta real dibujada entre las paradas y sus marcadores.
      \item Lista de \textit{Indicaciones de ruta} (en la captura, 20 tramos) calculada con la API de rutas de Google.
      \item Botón \textit{Ir con Google Maps} que abre el recorrido completo en la aplicación de mapas.
    \end{itemize}
    \vspace{0.2cm}

    \textbf{Lógica y Datos:}
    \begin{itemize}[leftmargin=*, itemsep=2pt]
      \item Si el GPS no responde, la ruta parte desde la primera parada; si la API no está disponible, se traza una línea directa entre paradas con tiempos estimados.
    \end{itemize}
  }

% --- PANTALLA 13: MapScreen ---
\pantallashowcase
  {Pantalla 13: Mapa Interactivo de Huecas}
  {lib/screens/map_screen.dart}
  {COMENSAL}
  {huecaOrange}
  {capturas/20_mapa.png}
  {17.5cm}
  {
    \proposito{Visualización geoespacial de las huecas de Quito con Google Maps.}

    \textbf{Componentes Clave:}
    \begin{itemize}[leftmargin=*, itemsep=2pt]
      \item Buscador por hueca, plato, sucursal o barrio.
      \item Filtro por zona (\textit{Todos, Centro, Norte, Sur, Valles}) con contador de puntos visibles (\textit{Mostrando todo: 20 puntos}).
      \item Marcadores de cada hueca y de sus sucursales, desde Cotocollao hasta Sangolquí.
      \item Botones de pantalla completa y de centrado en la ubicación actual.
    \end{itemize}
  }

% --- PANTALLA 14: ScanScreen ---
\pantallashowcase
  {Pantalla 14: Escáner de Códigos QR para Sellos}
  {lib/screens/scan_screen.dart}
  {COMENSAL}
  {huecaOrange}
  {capturas/21_scan.png}
  {17.5cm}
  {
    \proposito{Captura los códigos QR que muestra cada hueca para estampar sellos en la cartilla o desbloquear cromos secretos.}

    \textbf{Componentes Clave:}
    \begin{itemize}[leftmargin=*, itemsep=2pt]
      \item Visor de cámara en tiempo real con el paquete \texttt{mobile\_scanner}.
      \item Marco guía con línea de escaneo animada.
      \item Indicador superior \textit{``Apunta al QR del puesto''}.
      \item Controles de linterna y cambio de cámara.
    \end{itemize}
    \vspace{0.2cm}

    \textbf{Seguridad:}
    \begin{itemize}[leftmargin=*, itemsep=2pt]
      \item Cada código es de un solo uso: al escanearse se desactiva y la pantalla del dueño genera uno nuevo, lo que evita fraudes con capturas compartidas.
    \end{itemize}
  }

% --- PANTALLA 15: AlbumScreen ---
\pantalladoble
  {Pantalla 15: Mi Álbum --- Cartillas y Cromos por Zona}
  {lib/screens/album_screen.dart}
  {COMENSAL}
  {huecaOrange}
  {capturas/22_album.png}
  {capturas/25_cromos.png}
  {14cm}
  {
    \proposito{Pasaporte de gamificación del comensal: sellos acumulados por hueca y colección de cromos.}
    \begin{itemize}[leftmargin=*, itemsep=2pt]
      \item Tarjeta del usuario con nivel (\textit{Explorador}) y total de sellos.
      \item \textbf{Cartillas por Hueca (izquierda):} cada hueca visitada tiene su cartilla de sellos para canjear premios.
      \item \textbf{Cromos por Sector (derecha):} un álbum por zona (\textit{Centro, Norte, Sur, Valles}) con 5 cromos cada una y barra de progreso.
      \item Aviso para invitados: el progreso se guarda en el celular hasta que inicien sesión.
    \end{itemize}
  }

% --- PANTALLA 16: ProfileScreen ---
\pantalladoble
  {Pantalla 16: Perfil, Pasaporte y Accesibilidad}
  {lib/screens/profile_screen.dart}
  {COMENSAL}
  {huecaOrange}
  {capturas/23_perfil.png}
  {capturas/24_perfil_config.png}
  {14cm}
  {
    \proposito{Centro personal del comensal: progreso, favoritos, perfil gastronómico y configuración de accesibilidad.}
    \begin{itemize}[leftmargin=*, itemsep=2pt]
      \item Cabecera con avatar, nombre, correo y contadores: \textit{Huecas, Favoritos y Premios}.
      \item \textit{Mis Favoritos} y \textit{Mi Perfil Gastronómico} con el resultado del cuestionario y opción de repetirlo.
      \item \textbf{Accesibilidad e inclusión:} idioma (Español / English), tamaño de texto (A-, Normal, A+), tema claro u oscuro y \textbf{Alto Contraste}.
      \item \textit{Gestión de Cuenta}: tipo de sesión (Google, correo o modo invitado), tutorial de la app, \textit{¿Tienes una hueca? Sé parte de HueQuito}, ayuda, términos y privacidad, invitar a otros comensales, cerrar sesión y eliminar cuenta.
    \end{itemize}
  }

% --- PANTALLA 17: HuecaRequestScreen ---
\pantalladoble
  {Pantalla 17: Solicitud para Registrar una Hueca}
  {lib/screens/hueca_request_screen.dart}
  {COMENSAL}
  {huecaOrange}
  {capturas/26_solicitud_hueca.png}
  {capturas/26b_solicitud_form.png}
  {14cm}
  {
    \proposito{Permite que cualquier dueño de una hueca pida unirse a la plataforma; el administrador revisa y aprueba la solicitud.}
    \begin{itemize}[leftmargin=*, itemsep=2pt]
      \item Plan de inicio gratuito: \textit{100 sellos de cortesía}, sin costo de alta, y luego la suscripción mensual que fija el administrador (en la captura, \$4.99) con 7 días de gracia.
      \item \textbf{1. Datos de la Hueca:} nombre, zona de Quito, dirección y teléfono o WhatsApp.
      \item \textbf{2. Sazón y Especialidad:} plato estrella e historia de la hueca.
      \item \textbf{3. Contacto del Propietario:} nombre y correo de la futura cuenta de negocio, con aviso de que esa cuenta pasará a ser el panel del negocio.
      \item \textbf{4. Foto del Local o Plato} (opcional) para verificar más rápido.
    \end{itemize}
  }

% =============================================================================
% PARTE II: MÓDULO DE NEGOCIO (DUEÑO DE HUECA)
% =============================================================================
\newpage
\section{Módulo de Negocio (Dueño / Encargado de Hueca)}
\skipnewpagetrue

% --- PANTALLA 18: Guía del negocio + Dashboard ---
\pantalladoble
  {Pantalla 18: Guía de Inicio y Dashboard del Negocio}
  {lib/screens/business/dashboard_screen.dart}
  {DUEÑO / ENCARGADO}
  {huecaOrangeDark}
  {capturas/d00_guia_negocio.png}
  {capturas/d01_dashboard.png}
  {14cm}
  {
    \proposito{Panel de control con el que el dueño de la hueca (ejemplo: \textit{El Pretil}) mide el impacto de su local en la plataforma.}
    \begin{itemize}[leftmargin=*, itemsep=2pt]
      \item \textbf{Guía de tu Negocio (izquierda):} tutorial de 7 pasos que se abre la primera vez y resume las pestañas Perfil, Menú, Mi QR, Avisos y Dashboard, además del plan gratuito de 100 sellos.
      \item \textbf{Resumen de Impacto (derecha):} vistas totales del perfil en tiempo real vía Firestore.
      \item Tarjetas de métricas: \textit{Favoritos}, \textit{En Rutas} (tours de comensales que incluyen la hueca), \textit{Sellos Dados} y \textit{Reseñas}.
      \item Barra \textit{Fase Gratuita de Prueba} (avance hacia los 100 sellos) y barra inferior: \textit{Dashboard, Menú, Avisos, Mi QR, Perfil y Ver app}.
    \end{itemize}
  }

% --- PANTALLA 19: Menú digital por categorías ---
\pantalladoble
  {Pantalla 19: Menú Digital por Categorías y Asistente de IA}
  {lib/screens/business/menu_management_screen.dart}
  {DUEÑO / ENCARGADO}
  {huecaOrangeDark}
  {capturas/d02_menu.png}
  {capturas/d03_menu_anadir.png}
  {14cm}
  {
    \proposito{Gestión de los platos de la hueca, agrupados por categoría tal como los verá el comensal.}
    \begin{itemize}[leftmargin=*, itemsep=2pt]
      \item \textbf{Lista agrupada (izquierda):} encabezados por categoría (Desayunos, Bebidas, Postres, Platos Fuertes\dots), foto, nombre en español e inglés (badge \texttt{EN}), precio con dos decimales y botones de editar y eliminar.
      \item \textbf{Botón \faIcon{magic}\ ``Organizar por categorías con IA'':} Gemini clasifica todos los platos en las categorías del menú y los guarda ordenados.
      \item \textbf{Añadir al Menú (derecha):} \textit{Escanear Carta con Gemini IA (1 a 10 páginas)}, \textit{Escanear 1 Plato por Foto}, \textit{Subir Múltiples Fotos de Platos} o \textit{Añadir Plato Manualmente}.
      \item Las funciones de IA usan Gemini a través de Firebase AI Logic, sin claves de API dentro de la aplicación.
    \end{itemize}
  }

% --- PANTALLA 20: Formulario de plato y categorías ---
\pantalladoble
  {Pantalla 20: Formulario de Plato y Categorías Propias}
  {lib/screens/business/dish\_form\_screen.dart}
  {DUEÑO / ENCARGADO}
  {huecaOrangeDark}
  {capturas/d05_plato_categorias.png}
  {capturas/d06_nueva_categoria.png}
  {14cm}
  {
    \proposito{Alta y edición de cada plato, con la categoría que alimenta los filtros del menú del comensal.}
    \begin{itemize}[leftmargin=*, itemsep=2pt]
      \item Foto del plato, nombre y descripción en español e inglés con botón de \textit{Traducción Bilingüe Automática}, precio y disponibilidad en la carta.
      \item \textbf{Categorías sugeridas (izquierda):} Desayunos, Entradas, Sopas y Caldos, Platos Fuertes, Platos Típicos, Especialidades, Panadería, Postres, Bebidas y Otros.
      \item \textbf{\textit{+ Nueva categoría} (derecha):} el dueño crea sus propios grupos (por ejemplo \textit{Combos}), que aparecen como un chip más en el menú público.
    \end{itemize}
  }

% --- PANTALLA 21: Avisos ---
\pantalladoble
  {Pantalla 21: Promociones y Redactor con IA}
  {lib/screens/business/notification\_management\_screen.dart}
  {DUEÑO / ENCARGADO}
  {huecaOrangeDark}
  {capturas/d07_avisos.png}
  {capturas/d08_avisos_form.png}
  {14cm}
  {
    \proposito{Campañas y avisos que llegan al celular de los comensales suscritos a la hueca, previa aprobación del administrador.}
    \begin{itemize}[leftmargin=*, itemsep=2pt]
      \item \textbf{Menú del Día:} publica la sopa y el segundo del día de inmediato o programado.
      \item Avisos sobre el tiempo de aprobación (2 horas a 1 día) y el alcance de las notificaciones.
      \item Categorías: \textit{Menú del Día, Promoción, Noticia, Nuevo Plato y Evento}; sucursal destino.
      \item Título, mensaje con botón \textit{Ayuda con IA}, etiqueta opcional (\textit{2×1}) y salida \textit{Inmediata} o \textit{Programada}.
      \item \textbf{Vista previa en el celular (derecha):} muestra la notificación tal como la recibirá el cliente, y el botón \textit{Solicitar Autorización de Notificación}.
    \end{itemize}
  }

% --- PANTALLA 22: QR ---
\pantallashowcase
  {Pantalla 22: Emisor de QR Dinámico Antifraude}
  {lib/screens/business/qr_generator_screen.dart}
  {DUEÑO / ENCARGADO}
  {huecaOrangeDark}
  {capturas/d09_qr.png}
  {17.5cm}
  {
    \proposito{Entrega presencial de sellos de consumo con un código QR que no puede reutilizarse ni compartirse.}

    \textbf{Componentes Clave:}
    \begin{itemize}[leftmargin=*, itemsep=2pt]
      \item Sucursal activa (\textit{Matriz Principal}) y nombre de la hueca.
      \item Código QR generado en pantalla con el paquete \texttt{qr\_flutter}.
      \item \textbf{1 Solo Uso:} al escanearse, el código se desactiva y se regenera automáticamente.
      \item Selector \texttt{-} / \texttt{+} para indicar cuántos sellos se entregan según el consumo.
      \item Botones para regenerar el código y ver el historial de sellos emitidos.
    \end{itemize}
  }

% --- PANTALLA 23: Perfil del negocio ---
\pantalladoble
  {Pantalla 23: Perfil de la Hueca}
  {lib/screens/business/profile_management_screen.dart}
  {DUEÑO / ENCARGADO}
  {huecaOrangeDark}
  {capturas/d10_perfil.png}
  {capturas/d11_perfil_b.png}
  {14cm}
  {
    \proposito{Administración de los datos comerciales de la hueca que ven los comensales.}
    \begin{itemize}[leftmargin=*, itemsep=2pt]
      \item \textbf{Información General (izquierda):} portada, fotos de la hueca (máximo 5), locales y sucursales, nombre, descripción bilingüe, dirección escrita y punto exacto en el mapa (GPS).
      \item \textbf{Más opciones (derecha):} etiquetas, horario de atención, plato insignia (nombre, descripción e imagen).
      \item \textbf{Fidelización y Premios:} número de sellos necesarios para llenar la cartilla y premio por cartilla llena.
    \end{itemize}
  }

% --- PANTALLA 24: Sucursales ---
\pantalladoble
  {Pantalla 24: Locales y Sucursales}
  {lib/screens/business/branch\_management\_screen.dart}
  {DUEÑO / ENCARGADO}
  {huecaOrangeDark}
  {capturas/d12_sucursales.png}
  {capturas/d13_nuevo_local.png}
  {14cm}
  {
    \proposito{Red de locales de una misma hueca, cada uno con su ubicación y, opcionalmente, su encargado.}
    \begin{itemize}[leftmargin=*, itemsep=2pt]
      \item \textbf{Lista de locales (izquierda):} \textit{Matriz Principal} con dirección, zona y coordenadas GPS; indicador de \textit{Dueño General (Acceso Total)}.
      \item \textbf{Nuevo Local (derecha):} nombre, zona de Quito (Centro, Norte, Sur o Valles), dirección, punto en el mapa y teléfono.
      \item Encargado opcional por correo: esa persona recibe un enlace y administra solo su local.
    \end{itemize}
  }


% =============================================================================
% PARTE III: MÓDULO DE ADMINISTRACIÓN (SUPERADMIN)
% =============================================================================
\newpage
\section{Módulo de Administración (Superadministrador)}
\skipnewpagetrue

% --- PANTALLA 25: Acceso y Centro de Control ---
\pantalladoble
  {Pantalla 25: Acceso al Panel y Centro de Control}
  {lib/screens/admin/admin_dashboard_screen.dart}
  {SUPERADMIN}
  {adminNavy}
  {capturas/a00_perfil_admin.png}
  {capturas/a01_dashboard.png}
  {14cm}
  {
    \proposito{Tablero maestro para monitorear el estado global de la plataforma y entrar a cada bandeja de gestión.}
    \begin{itemize}[leftmargin=*, itemsep=2pt]
      \item \textbf{Acceso (izquierda):} la tarjeta \textit{Panel de Administración} aparece en el perfil únicamente para la cuenta de superadministrador.
      \item \textbf{Métricas clave (derecha):} \textit{Huecas} (20, todas activas), \textit{Usuarios}, \textit{Avisos Pendientes}, \textit{Sectores y Cromos} (4 zonas), \textit{Solicitudes Huecas}, \textit{Bandeja Soporte}, \textit{Suscripciones} y \textit{Quiz y ADN}; cada tarjeta abre su pantalla.
      \item Barra superior con distintivo \textbf{SUPERADMIN} y salida rápida; barra inferior: \textit{Dashboard, Huecas, Avisos, Pagos y Usuarios}.
    \end{itemize}
  }

% --- PANTALLA 26: Directorio de huecas + edición ---
\pantalladoble
  {Pantalla 26: Directorio y Edición de Huecas}
  {lib/screens/admin/admin_huecas_screen.dart}
  {SUPERADMIN}
  {adminNavy}
  {capturas/a02_huecas.png}
  {capturas/a03_editar_hueca.png}
  {14cm}
  {
    \proposito{Gestión centralizada de las huecas: visibilidad, datos, ubicación y dueño asignado.}
    \begin{itemize}[leftmargin=*, itemsep=2pt]
      \item \textbf{Directorio (izquierda):} buscador por nombre, plato, menú o zona; filtros \textit{Todas, Activas, Pausadas, Sin Dueño}; acceso a \textit{Sectores y Cromos}.
      \item Tarjetas con foto, zona, calificación, estado del dueño (\textit{Dueño Asignado} o \textit{Sin Dueño}), número de sucursales, interruptor de visibilidad y acciones \textit{Locales, Dueño, Editar y Eliminar}.
      \item \textbf{Editar Hueca (derecha):} estado, nombre, zona (Centro, Norte, Sur o Valles), rango de precio, teléfono, dirección con selector en mapa y coordenadas, y cuenta de negocio vinculada con buscador para cambiar de dueño.
      \item Regla visible: cada cuenta de negocio solo puede estar asignada a una hueca.
    \end{itemize}
  }

% --- PANTALLA 27: Sectores y cromos ---
\pantallashowcase
  {Pantalla 27: Sectores y Cromos del Álbum}
  {lib/screens/admin/admin_sectors_screen.dart}
  {SUPERADMIN}
  {adminNavy}
  {capturas/a04_sectores.png}
  {17.5cm}
  {
    \proposito{Organiza las huecas en las cuatro zonas del álbum coleccionable y valida que cada una tenga su cromo secreto.}

    \textbf{Componentes Clave:}
    \begin{itemize}[leftmargin=*, itemsep=2pt]
      \item Resumen: 4 sectores, huecas con zona confirmada (20/20) y huecas con cromo (20/20).
      \item Buscador, botón \textit{Confirmar Todo} y filtros \textit{Todos, Por Confirmar, Sin Cromo Secreto}.
      \item Grupos por zona con estado \textit{Validado} y número de cromos listos.
      \item Por hueca: zona confirmada con opción \textit{Cambiar}, nombre del cromo (por defecto \textit{Cromo de} y el nombre de la hueca), QR listo para escanear y \textit{Editar Cromo}.
    \end{itemize}
  }

% --- PANTALLA 28: Moderación de avisos ---
\pantallashowcase
  {Pantalla 28: Moderación de Avisos y Notificaciones}
  {lib/screens/admin/admin_notifications_screen.dart}
  {SUPERADMIN}
  {adminNavy}
  {capturas/a05_avisos.png}
  {17.5cm}
  {
    \proposito{Control editorial de las promociones que envían los dueños antes de que lleguen a los comensales.}

    \textbf{Componentes Clave:}
    \begin{itemize}[leftmargin=*, itemsep=2pt]
      \item Pestañas \textit{Pendientes} (con contador) e \textit{Historial Completo}.
      \item Accesos a otras bandejas: \textit{Solicitudes} de huecas nuevas y \textit{Soporte} de negocios.
      \item Tarjeta del aviso con hueca emisora, categoría, etiqueta (\textit{2×1}) y tipo de emisión.
      \item Simulación de cómo saldrá la notificación en Android.
      \item Botones \textit{Rechazar} y \textit{Aprobar y Emitir}.
    \end{itemize}
  }

% --- PANTALLA 29: Solicitudes de huecas ---
\pantallashowcase
  {Pantalla 29: Solicitudes de Nuevas Huecas}
  {lib/screens/admin/admin_solicitudes_screen.dart}
  {SUPERADMIN}
  {adminNavy}
  {capturas/a06_solicitudes.png}
  {17.5cm}
  {
    \proposito{Revisión de las solicitudes enviadas desde \textit{¿Tienes una hueca? Sé parte de HueQuito}.}

    \textbf{Componentes Clave:}
    \begin{itemize}[leftmargin=*, itemsep=2pt]
      \item Pestañas \textit{Pendientes, Aprobadas, Rechazadas y Todas}, con buscador por nombre, zona o dueño.
      \item Tarjeta con nombre de la hueca, zona, fecha, dueño, correo, teléfono, dirección, plato estrella e historia.
      \item \textit{Aprobar} crea la hueca y la cuenta de negocio; \textit{Rechazar} registra el motivo.
    \end{itemize}
  }

% --- PANTALLA 30: Suscripciones y pagos ---
\pantallashowcase
  {Pantalla 30: Suscripciones y Pagos}
  {lib/screens/admin/admin_subscriptions_screen.dart}
  {SUPERADMIN}
  {adminNavy}
  {capturas/a07_pagos.png}
  {17.5cm}
  {
    \proposito{Seguimiento del modelo de negocio: 100 sellos gratuitos por hueca y luego una mensualidad.}

    \textbf{Componentes Clave:}
    \begin{itemize}[leftmargin=*, itemsep=2pt]
      \item Gráfica \textit{Próximos a Cobro} con la meta (100 sellos) y el cobro mensual configurado.
      \item Contadores \textit{Cerca a Cobro, Meta Alcanzada y En Prueba}, y avance de cada hueca hacia los 100 sellos.
      \item \textit{Tarifas y Datos Bancarios}: precio mensual, meta de sellos, días de gracia y cuenta para transferencias; este precio es el que ven los dueños en la solicitud y en su guía.
      \item Pestañas de comprobantes: \textit{Por Validar, Al Día y En Gracia / Corte}.
    \end{itemize}
  }

% --- PANTALLA 31: Usuarios + alta de negocio ---
\pantalladoble
  {Pantalla 31: Usuarios y Alta Segura de Negocios}
  {lib/screens/admin/admin_users_screen.dart}
  {SUPERADMIN}
  {adminNavy}
  {capturas/a08_usuarios.png}
  {capturas/a10_crear_cuenta.png}
  {14cm}
  {
    \proposito{Administración de cuentas de comensales, negocios y administración.}
    \begin{itemize}[leftmargin=*, itemsep=2pt]
      \item \textbf{Directorio (izquierda):} buscador por nombre, correo, hueca o sucursal; filtros \textit{Todos, Con Hueca / Local, Sucursales, Sin Asignar}; insignias \texttt{ADMIN} y \texttt{NEGOCIO} con la hueca vinculada; avatar con iniciales si la cuenta no tiene foto.
      \item Acciones por usuario: editar rol, eliminar y \textit{Reenviar enlace de contraseña}.
      \item \textbf{Crear Cuenta de Negocio (derecha):} nombre y correo del dueño; no se escribe ninguna contraseña: el dueño recibe un correo para definirla. Selector de hueca con la regla de una hueca por cuenta.
    \end{itemize}
  }

% --- PANTALLA 32: Quiz ---
\pantallashowcase
  {Pantalla 32: Configuración del Quiz Gastronómico}
  {lib/screens/admin/admin_quiz_screen.dart}
  {SUPERADMIN}
  {adminNavy}
  {capturas/a09_quiz.png}
  {17.5cm}
  {
    \proposito{Editor de las preguntas del cuestionario inicial y de su relación con los perfiles gastronómicos.}

    \textbf{Componentes Clave:}
    \begin{itemize}[leftmargin=*, itemsep=2pt]
      \item Explicación de la asignación automática: cada opción suma un punto a un perfil (\textit{Sopero, Carnívoro, Callejero, Dulcero, Picador, Aventurero}) y gana el de mayor puntaje.
      \item Lista de preguntas (4) con texto en español e inglés, fotografía y opciones con su perfil asignado.
      \item Botones para crear, editar y eliminar preguntas.
    \end{itemize}
  }


\end{document}
